% HD-N20. GPT-6 Astra Pro /root/imprimitive. 21 November 2026.
% Inherited reasoning effort; the numerical setting is not exposed.
\section{HD-N20: an actual semiregular product-surface ideal and its whole-character orbit}
\label{sec:product-surface-ideal}

\paragraph{The object and the two conclusions.}
Let $C_1,C_2$ be smooth nonhyperelliptic curves of genus three,
let $j_i:C_i\hookrightarrow J_i=J(C_i)$ be their Abel--Jacobi
embeddings, and put
\[
 X=J_1\times J_2,\qquad S=C_1\times C_2\subset X,\qquad F=I_S.
                                                               \tag{N20.1}
\]
This is an actual coherent sheaf on a smooth sixfold and hence a
perfect complex with negative self-Ext zero. Every categorical trace
$\sigma_F^k$ is injective, and its self-Ext dimensions are
\[
                   (1,12,54,94,54,12,1).                       \tag{N20.2}
\]
The proof below checks the maps in every degree, not only their
Euler characteristic.

No whole character in the actual Fourier--Mukai, isogeny pullback or
pushforward, arbitrary line-bundle twist, translation, shift and
derived-dual orbit of $F$, even multiplied by a nonzero rational
scalar, is
\[
 e^B\bigl(P(\theta)+\lambda t\bigr),\qquad
 \deg P\le6,\quad 0\ne\lambda t\in W_K(A),                       \tag{N20.3}
\]
for a polarized Weil sixfold of signature$(3,3)$, where
$B\in H^{1,1}(A)\cap H^2(A,\mathbf R)$ and $P$ has real
coefficients. This includes rational algebraic $B$.
Here the orbit contains those specified actual functors, with
arbitrary mixing of the two original factors. Arbitrary other
retracts or extensions are not substituted for those functors.
The Fermat-quartic construction at the end gives an actual object
on $E_i^6$ with a coherent translation group of order64 and the
same invariant self-Ext and trace assertions.

\paragraph{1. The exact character map.}
Write $\theta_i$ for the principal Jacobian polarization and also
for its pullback to $X$. The Abel--Jacobi map identifies
$H^1(J_i)$ with $H^1(C_i)$ and its alternating intersection pairing.
Evaluating on a symplectic basis therefore gives the Poincare formula
$j_{i,*}[C_i]=\theta_i^2/2$: pairing both sides with each product
of two degree-one basis classes gives the same alternating pairing.
Grothendieck--Riemann--Roch for the actual embedding $j_i$ gives
\[
 a_i:=\operatorname{ch}(j_{i,*}\mathcal O_{C_i})
       =\frac{\theta_i^2}{2}-2[\mathrm{pt}]
       =\frac{\theta_i^2}{2}-\frac{\theta_i^3}{3}.               \tag{N20.4}
\]
Indeed the Todd class of $J_i$ is1 and
$\int_{C_i}\operatorname{td}_1(C_i)=1-g(C_i)=-2$.
The published GRR statement is Fulton, \emph{Intersection Theory},
second edition, Theorem15.2; the publisher's chapter statement
and formula are at
\href{https://doi.org/10.1007/978-1-4612-1700-8_16}
{DOI10.1007/978-1-4612-1700-8\_16}, pp.280--304.

The actual closed product embedding gives
$\mathcal O_S=j_{1,*}\mathcal O_{C_1}\boxtimes
j_{2,*}\mathcal O_{C_2}$. Its character is $a_1a_2$ by the
external-product morphism in $K$-theory. The short exact sequence
$0\to I_S\to\mathcal O_X\to\mathcal O_S\to0$ consequently gives
the whole character
\[
                         v=\operatorname{ch}(F)=1-a_1a_2.     \tag{N20.5}
\]
These maps specify the product-surface ideal itself; no tensor
product of the two curve ideals is used in(N20.5).

\paragraph{2. The actual Hochschild maps in all degrees.}
For each threefold choose invariant forms
$z_{ij}\in H^{1,0}(J_i)$ and $b_{ij}\in H^{0,1}(J_i)$ with
$\theta_i=\sum_{j=1}^3 z_{ij}b_{ij}$. Put
\[
 U_i=H^{0,1}(J_i)\oplus H^{1,0}(J_i)^*,\quad
 W_{ij}=b_{ij}\wedge,\quad D_{ij}=\iota_{z_{ij}}.
\]
The action on exterior forms is literal wedge and contraction.
The published trace/evaluation identity is Buchweitz--Flenner,
Adv. Math.217 (2008),243--281, Proposition6.2.1,
Corollary6.3.2 and Theorem6.4.2, author version
\href{https://arxiv.org/pdf/math/0606730}{arXiv:math/0606730},
pp.29--31. Since the Todd class is1, it identifies this action
on$v$ with $c_F^k=\sigma_F^k\operatorname{ev}_F^k$.

Here are the complete one-factor maps needed for the calculation.
Suppress $i$ and put $a=\theta^2/2-\theta^3/3$ and
$T=H^{0,1}(J)$. The degree-one map $U\to H^*(J)$ acting on$a$
is injective, of rank6:
\[
 W_j a=\tfrac12\theta^2 b_j,\qquad
 D_j a=(\theta-\theta^2)b_j.                                  \tag{N20.6}
\]
The three degree-three lowest components $\theta b_j$ are
independent, as are the three degree-five forms $\theta^2b_j$.
Their distinct bidegrees prove injectivity. The map acting on1
has image$T$, in degree one.

In degree two, for $j<k$ and the indicated operator ordering,
\[
 \begin{array}{c|c|c}
        &\text{action on }1&\text{action on }a\\ \hline
 W_jW_k&b_jb_k&0\\
 W_jD_k&0&\theta b_jb_k\\
 D_jD_k&0&(1-2\theta)b_jb_k
 \end{array}                                                  \tag{N20.7}
\]
All other mixed columns have the corresponding exterior signs
or are zero. Thus the action on$a$ has image
\[
                   \bigwedge^2T\oplus\theta\bigwedge^2T
\]
of dimension6. The joint action on1 and$a$ maps onto the
nine-dimensional direct sum
\[
 \bigwedge^2T\ \oplus\
       \bigl(\bigwedge^2T\oplus\theta\bigwedge^2T\bigr).          \tag{N20.8}
\]
In degree three both individual images are the line$\bigwedge^3T$,
and their joint map has rank2: $W_1W_2W_3$ acts by
$(b_1b_2b_3,0)$, whereas $W_1D_2D_3$ acts by
$(0,b_1b_2b_3)$. In higher degrees the action on1 and on$a$
vanishes, since the antiholomorphic degree minus the holomorphic
degree would exceed three.

The Hochschild--Kunneth decomposition is
\[
 HH^k(X)=\bigoplus_{r+s=k}\bigwedge^rU_1\otimes\bigwedge^sU_2.
\]
On the $(r,s)$ summand the map acting on(N20.5) is the difference
of the tensor products of the two actions on1 and the two actions
on$a_i$, with the common Koszul sign. All outputs in that summand
have Hodge differences $(q_1-p_1,q_2-p_2)=(r,s)$.
Thus different summands have disjoint images. The complete table is
\[
 \begin{array}{c|rrrr}
 r\backslash s&0&1&2&3\\ \hline
 0&1&6&9&2\\
 1&6&36&45&9\\
 2&9&45&36&6\\
 3&2&9&6&1
 \end{array}                                                  \tag{N20.9}
\]
We justify every entry. The $(0,s)$ and $(r,0)$ entries follow
from(N20.6)--(N20.8) and independence of1 and the other$a_i$.
For $(1,1)$, the tensor of the two maps acting on$a_i$ is
injective of rank36. Its outputs have degree at least six,
whereas the additional tensor of the two actions on1 has
degree two. This retains rank36.

For $(1,2)$ the product of the two actions on$a_i$ has rank
$6\cdot6=36$ and degree at least five. The additional term
acting on1 has degree three. On the kernel of the former map,
its image has dimension9: the three $W_jW_k$ columns in the
second factor lie in that kernel and map onto$\bigwedge^2T_2$,
while the first action on1 has rank3. There cannot be more
than those nine extra dimensions. This gives45, and the
$(2,1)$ calculation is the exchanged one.

For $(1,3)$ the degree-three joint map is onto its two lines,
so the image is the sum of the first factor's rank-three
action on1 and rank-six action on$a_1$. Their degrees are
disjoint, giving9; the exchanged case gives9 as well.
For $(2,2)$, the product of the actions on$a_i$ already maps
onto
\[
 (\bigwedge^2T_1\oplus\theta_1\bigwedge^2T_1)
 \otimes
 (\bigwedge^2T_2\oplus\theta_2\bigwedge^2T_2),
\]
a space of dimension36. The additional product of the
actions on1 lies in the first tensor summand of that space.
The surjectivity of both joint maps(N20.8) permits the
$a_i$ inputs to be chosen with their action on1 equal to zero,
so cancellation does not shrink this image.
For $(2,3)$ the image is the sum of the action on1 and on$a_1$;
the first is contained in the second, giving6. The exchanged
case gives6. Finally $(3,3)$ is the single top antiholomorphic
line, reached by the all-wedge operator. This proves the table.

Summing its diagonals gives the exact ranks of the actual maps:
\[
        \bigl(\operatorname{rank}c_F^k\bigr)_{k=0}^6
                         =(1,12,54,94,54,12,1).                \tag{N20.10}
\]
In particular the degree-two rank is $9+36+9=54$.

\paragraph{3. The local ideal resolution and the Ext bounds.}
Let $N_i=N_{C_i/J_i}$. The normal exact sequence gives
$h^0(N_i)=3$: the additional tangent directions would be the
kernel of $H^1(T_{C_i})\to H^1(T_{J_i}|_{C_i})$.
Its dual contains
\[
 \operatorname{Sym}^2H^0(K_{C_i})\longrightarrow H^0(K_{C_i}^2),
\]
an isomorphism for a plane quartic, since both dimensions are6
and its canonical plane model has no quadratic equation.
Also $\det N_i=K_{C_i}$ and $\chi(N_i)=4+2(1-3)=0$.
Thus $h^1(N_i)=3$. These are the normal maps already used
for the actual curve ideal in HD-N02.

The normal bundle of the actual product embedding is
$N=N_{S/X}=p_1^*N_1\oplus p_2^*N_2$.
The ordinary coherent Kunneth maps give
\[
 h^0(N)=6,\qquad h^1(N)=2(3+3\cdot3)=24,\qquad
 h^0(\bigwedge^2N)=3+3\cdot3+3=15.                             \tag{N20.11}
\]
The embedding is regular of codimension four. A local Koszul
resolution of its ideal, ending in
$\bigwedge^2N^*\to N^*\to I_S$, has zero differentials after
applying $\mathcal Hom(-,\mathcal O_S)$. Consequently its actual
local Ext sheaves are
\[
 \mathcal Ext_X^q(I_S,\mathcal O_S)=\bigwedge^{q+1}N
                      \quad(0\le q\le3).                      \tag{N20.12}
\]
The local-to-global spectral sequence therefore gives the
bounds, without assuming its degeneration,
\[
 \operatorname{hom}(I_S,\mathcal O_S)=6,\qquad
 \operatorname{ext}^1(I_S,\mathcal O_S)
              \le h^1(N)+h^0(\bigwedge^2N)=39.                 \tag{N20.13}
\]
The dual Koszul resolution gives
$R\mathcal Hom_X(\mathcal O_S,\mathcal O_X)
 =j_*\det N[-4]$.
Dualizing the ideal triangle thus identifies the local derived
dual of$I_S$ with$\mathcal O_X$ in degrees at most two.
After global cohomology this gives the actual identifications
\[
 \operatorname{Ext}^1(I_S,\mathcal O_X)=H^1(\mathcal O_X),
 \qquad
 \operatorname{Ext}^2(I_S,\mathcal O_X)=H^2(\mathcal O_X),        \tag{N20.14}
\]
of dimensions6 and15.

Now apply $R\operatorname{Hom}(I_S,-)$ to the specified ideal
triangle. Its actual connecting maps give the exact segments
\[
 \begin{split}
 \operatorname{Hom}(I_S,\mathcal O_S)&\longrightarrow
 \operatorname{Ext}^1(I_S,I_S)\longrightarrow
 \operatorname{Ext}^1(I_S,\mathcal O_X),\\
 \operatorname{Ext}^1(I_S,\mathcal O_S)&\longrightarrow
 \operatorname{Ext}^2(I_S,I_S)\longrightarrow
 \operatorname{Ext}^2(I_S,\mathcal O_X).
 \end{split}
\]
Equations(N20.13)--(N20.14) yield
\[
              \operatorname{ext}^1(F,F)\le12,\qquad
              \operatorname{ext}^2(F,F)\le54.                  \tag{N20.15}
\]
No connecting map has been declared zero. The actual character
ranks12 and54 in(N20.10) give the reverse inequalities and
injectivity of$\sigma_F^1,\sigma_F^2$. In particular this proves
semiregularity before any use of higher trace degrees.

\paragraph{4. The remaining degrees and the exact Euler calculation.}
An endomorphism of the ideal restricts to a function on$X\setminus S$.
It extends to$X$ since the complement has codimension four and$X$
is smooth. Thus $\operatorname{End}(F)=\mathbb C$.
Serre duality now gives the dimensions in degrees4,5,6 as54,12,1.
For the remaining middle dimension use the exact RR calculation.
The first nonconstant term in$v$ has Chern degree four, so
the product of any two nonconstant terms of$v^\vee$ and$v$
vanishes in dimension six. Duality has sign$+1$ in Chern
degree six, and
\[
 \int_Xa_1a_2=\chi(\mathcal O_{C_1})\chi(\mathcal O_{C_2})=4.
\]
It follows that
\[
                   \chi(F,F)=\int_Xv^\vee v=-8.               \tag{N20.16}
\]
The odd Chern-degree-five part has its dual sign reversed;
it contributes nothing to the integral. Thus no identification
$v^\vee=v$ is being made.
Equation(N20.16) gives
$-8=2-24+108-\operatorname{ext}^3(F,F)$, hence
$\operatorname{ext}^3(F,F)=94$. Negative self-Ext vanishes because
$F$ is a sheaf, and degrees above six vanish by Serre duality.

The complete rank table(N20.10) now equals these Ext dimensions
in every degree. Since $c_F^k=\sigma_F^k\operatorname{ev}_F^k$,
every$\sigma_F^k$ is injective. This proves(N20.2) and the
all-degree trace assertion through the explicit maps just checked.

\paragraph{5. The whole-character orbit obstruction.}
The exact balanced-sixfold calculation HD-R06 gives
\[
 \operatorname{rank}c^2_{\,e^B(P(\theta)+\lambda t)}
       =30+9\,\operatorname{rank}M_P
                     \in\{30,39,48,57\},                       \tag{N20.17}
\]
where $M_P=(m_{i+j})_{0\le i\le4,\,0\le j\le2}$ and
$m_j=j![z^j]P$. Its two Weil blocks have distinct weights
and ranks15 each; the mixed block consists of nine copies
of this displayed matrix. This statement uses a nonzero
rational Weil class, whose two conjugate components are
both nonzero.

We retain the exact maps which make rank an orbit invariant.
An actual Fourier--Mukai equivalence has its real Clifford
intertwiner $L c(w)L^{-1}=c(fw)$, with$f$ carrying the Hodge
subspace$U=H^{0,1}\oplus H^{1,0*}$ isomorphically to the
target one. The actual graph-kernel/point-functor map and the
Poincare swap identifying this$U$ are those proved in
HD-G19 and HD-N18, using Golyshev--Lunts--Orlov,
Proposition4.3.7 and Lemma4.3.7.1, and Orlov's graph-kernel
description. Thus the degree-two action diagrams commute
through$\bigwedge^2f$ and$L$.

For an isogeny with cohomology pullback$T$ the map is
$(T,(T^{-1})^*)$ on$H^1\oplus H^{1*}$; proper pushforward
on characters is its inverse multiplied by the degree, by
GRR and$p_*p^*=\deg p$. A line twist with class$D$ has the
literal intertwiner
\[
 e^D c(v+\xi)e^{-D}=c(v-\iota_\xi D+\xi);
\]
its$(1,1)$ condition preserves$U$, including mixed line
classes. Translation fixes cohomology and a shift multiplies
the character by a sign. Derived duality acts in degree$2j$
by$(-1)^j$; on degree-two Clifford products it conjugates
to$-c(f_{\mathrm d}u)c(f_{\mathrm d}v)$, where
$f_{\mathrm d}(x,\xi)=(x,-\xi)$ is the actual anti-isometry
from HD-N18. It too preserves$U$ and rank. Nonzero rational
multiplication of a character preserves rank as well.
Every listed operation therefore preserves the source
rank54 from(N20.10), which is absent from(N20.17).
This proves the entire stated actual-orbit exclusion.
It is independent of any subsequently chosen coherent action.

\paragraph{6. The actual CM object and its coherent action.}
Take both curves to be the Fermat plane quartic. HD-N01 constructs
the actual degree-eight isogeny
\[
 q:E_i^3\longrightarrow J(C),\qquad
 q(x_1,x_2,x_3)=\sum_{\nu=1}^3 f_\nu^*(x_\nu),\qquad
                         q^*\theta=2H.
\]
Here$f_\nu$ are the three explicitly written elliptic quotient
maps, and the diagonal norm maps are$[2]$ with all cross
norm maps zero. Thus the isogeny and its polarization scalar
have already been proved, not inferred solely from CM type.

Put$p=q\times q$ and
\[
 \widetilde F=p^*I_{C\times C}=I_{p^{-1}(C\times C)}
                  \quad\text{on }E_i^6.                       \tag{N20.18}
\]
Flat pullback gives the displayed actual ideal equality.
The group$G_0=\ker p$ has order64 and its translation action
on$\widetilde F$ has the canonical coherent descent structure.
The exact descent maps identify
\[
 \operatorname{Ext}^k(\widetilde F,\widetilde F)^{G_0}
             \simeq\operatorname{Ext}^k(F,F).
\]
Isogeny pullback on Hochschild homology is an isomorphism;
naturality of trace therefore gives injectivity in every
degree on these invariant Ext groups. This is in particular
$G_0$-semiregularity and hence weak$G_0$-semiregularity in
Perry Definition2.6 and Lemma2.7. Its whole character is
\[
 1-\left(2H_1^2-\tfrac83H_1^3\right)
     \left(2H_2^2-\tfrac83H_2^3\right).                         \tag{N20.19}
\]
The invariant dimensions are precisely(N20.2); no equality
with the entire nonequivariant self-Ext is asserted.
The rank54 obstruction remains valid on this CM source by
the actual isogeny map above.

\paragraph{Check record.}
The companion \texttt{product\_surface\_ideal\_certificate.py}
rebuilds the character, all66 degree-two operators and their
three Kunneth blocks, and supplies nonzero exact integer
minors of sizes9,36,9. It also reconstructs all sixteen
higher Kunneth blocks and checks the complete rank vector
at two finite fields. Only disclosed elementary wedge and
contraction utility definitions are reused from HD-N19;
no matrices or stored ranks are imported. The universal
rank proof is(N20.6)--(N20.10); the local Ext and coherent
descent arguments are separate from these finite checks.

